\section{Introduction}
\label{sec:introduction}

Indonesia's Ministry of Health Regulation No.\ 24/2022 mandates that every
healthcare facility operate Electronic Health Record (EHR) systems that
guarantee the confidentiality, integrity, and availability of patient
data~\cite{permenkes24_2022}. These guarantees extend beyond access control:
a credible audit trail---one that can reconstruct who accessed or modified a
record, when, and from which component---is a prerequisite for regulatory
compliance, forensic investigation, and insider-threat
detection~\cite{cert2022insiderthreat}.

DecMed addressed the access-control dimension of EHR security by combining
Capability-Based Access Control (CapBAC) enforced via IOTA smart contracts,
Proxy Re-Encryption (PRE) for key delegation without exposing plaintext, and
IPFS for off-chain clinical data storage~\cite{DecMed}. Despite these
contributions, DecMed's audit capabilities exhibit three structural
weaknesses. First, on-chain capability entries record \emph{granted
permissions}, not \emph{realized accesses}; the system cannot reconstruct who
actually read a record at a given time. Second, operations such as access
revocation overwrite existing state rather than appending a new entry,
violating the append-only invariant required for non-repudiation. Third,
logging is scattered across independent components with heterogeneous formats,
no central collector, and no retention strategy against IOTA's data-pruning
mechanism.

This paper makes two contributions:
\begin{enumerate}
    \item A systematic derivation of audit event requirements for DecMed using
    STRIDE threat modeling applied per-interaction on a Data Flow Diagram
    (DFD), yielding 13 event types that cover all six STRIDE categories.

    \item The design and implementation of the \textbf{Audit Log Service
    (ALS)}, a standalone Rust service that enforces tamper-evidence through
    SHA-256 hash-chaining, delivers long-term immutability through IPFS
    pinning and IOTA Move frozen objects, and integrates non-intrusively with
    DecMed's existing infrastructure.
\end{enumerate}

The remainder of this paper is organized as follows. Section~\ref{sec:background}
reviews EHR accountability standards, the DecMed framework, and relevant
security primitives. Section~\ref{sec:threat-analysis} presents the STRIDE
threat model and derives the audit event catalog. Section~\ref{sec:design}
describes the ALS architecture and key design decisions.
Section~\ref{sec:implementation} covers implementation details.
Section~\ref{sec:evaluation} evaluates correctness and security properties.
Section~\ref{sec:conclusion} concludes and outlines future work.